\chapter{Comparativa de modelos candidatos}
\label{cap:AnexoA}

Se consideraron diez modelos preentrenados como posible \emph{encoder} de texto para el clasificador CIE-10, seis de ellos evaluados experimentalmente. Los criterios de selección fueron: rendimiento en el corpus CodiEsp, especialización en español clínico, longitud de ventana de contexto y coste computacional de inferencia y entrenamiento.

\section{Modelos evaluados}

\subsection{\texttt{PlanTL-GOB-ES/roberta-base-biomedical-clinical-es}}

Modelo RoBERTa~\cite{liu2019roberta} de 125M parámetros entrenado por el Barcelona Supercomputing Center sobre corpus clínico en español: historias clínicas, informes radiológicos y notas médicas. Ventana de contexto de 512 tokens. Ofrece el mejor equilibrio entre especialización médica y velocidad de inferencia en CPU. Obtuvo un F1 de 0{,}76 en las pruebas de nivel~1 (clasificación de capítulos).

\subsection{\texttt{PlanTL-GOB-ES/roberta-large-bne}}

Variante de 355M parámetros del mismo grupo~\cite{liu2019roberta} (deprecado; se recomienda migrar a modelos BSC-LT). Mayor capacidad pero requiere GPU con al menos 16\,GB de VRAM para un entrenamiento viable. No se evaluó experimentalmente por restricciones de cómputo; se descartó para los niveles en los que el coste no estaba justificado.

\subsection{\texttt{PlanTL-GOB-ES/longformer-base-4096-biomedical-clinical-es}}

Variante del modelo clínico del Barcelona Supercomputing Center con ventana de contexto ampliada a 4.096 tokens mediante atención dispersa (\emph{Longformer}~\cite{beltagy2020longformer}), pensada para procesar informes largos sin segmentación. Se entrenó experimentalmente, pero se descartó por su elevado coste computacional ($\sim$70\,s por época frente a los $\sim$11\,s de RigoBERTa-Clinical, unas 6,4$\times$ más lento) sin mejora de rendimiento (F1-micro de 0{,}416 sobre los 809 códigos de bloque).

\subsection{\texttt{IIC/RigoBERTa-2.0}}

Modelo de propósito general basado en XLM-RoBERTa-large ($\sim$600M parámetros), desarrollado por el Instituto de Ingeniería del Conocimiento~\cite{vacaserrano2022rigoberta} sobre corpus español de gran escala. Ventana de contexto de 512 tokens (secuencias más largas se procesan mediante segmentación con solapamiento). Obtuvo un F1 de 0{,}74 en nivel~1 y se evaluó como base generalista frente a su variante clínica.

\subsection{\texttt{IIC/RigoBERTa-Clinical}}

Variante clínica de RigoBERTa-2.0 desarrollada por el Instituto de Ingeniería del Conocimiento~\cite{garciasubies2026rigobertaclinical}, especializada en textos médicos en español mediante entrenamiento adicional sobre corpus clínico. Comparte la arquitectura XLM-RoBERTa-large de la versión base. Obtuvo el mejor F1 experimental (0{,}799 en nivel~1) y fue seleccionada como modelo base para el entrenamiento final en producción.

\subsection{\texttt{jhu-clsp/mmBERT-base}}

Encoder multilingüe moderno (\emph{Modern Multilingual BERT}) desarrollado por el Centro de Procesamiento del Lenguaje y del Habla de la Universidad Johns Hopkins~\cite{marone2025mmbert}, basado en la arquitectura ModernBERT. Ventana de contexto de 8.192 tokens; entrenado sobre más de 1800 lenguas. Obtuvo un F1-micro de 0{,}12 sobre los 809 códigos de bloque por ausencia de especialización en español clínico y fue descartado.

\subsection{\texttt{answerdotai/ModernBERT-base}}

Arquitectura de nueva generación (2024) con mecanismo Flash Attention y ventana de contexto de 8.192 tokens~\cite{warner2024modernbert}. Obtuvo un F1-micro de 0{,}601 en las pruebas de nivel~1, inferior a RigoBERTa-Clinical, principalmente por operar en inglés frente a un corpus en español. Queda identificado como candidato para iteraciones futuras.

\subsection{\texttt{dccuchile/bert-base-spanish-wwm-cased} (BETO)}

Modelo BERT~\cite{devlin2019} de 110M parámetros entrenado sobre corpus en español de propósito general~\cite{canete2020beto}. Sin especialización médica. Descartado sin evaluación experimental propia.

\subsection{\texttt{sentence-transformers/paraphrase-multilingual-MiniLM-L12-v2}}

Modelo de 33M parámetros optimizado para generación de embeddings de similitud semántica~\cite{reimers2019sentencebert}. No es adecuado para clasificación directa; se evaluó únicamente como posible componente de una arquitectura basada en recuperación (\emph{RAG}).

\subsection{\texttt{distilbert-base-multilingual-cased}}

Versión destilada de BERT multilingüe con 66M parámetros~\cite{sanh2019distilbert}. Un 40\,\% más rápido que BERT base con pérdida de precisión moderada. Sin especialización en español médico; considerado como alternativa para entornos con recursos de cómputo muy limitados.

\section{Tabla resumen}

\begin{table}[h]
\centering
\caption{Comparativa de los modelos candidatos evaluados.}
\label{tab:modelos}
\begin{tabular}{lrccl}
\hline
\textbf{Modelo} & \textbf{Params} & \textbf{Ventana} & \textbf{F1 (tarea)} & \textbf{Dominio} \\
\hline
RoBERTa-BSC (base) & 125M & 512 & 0{,}778 & Clínico ES \\
RoBERTa-BSC (large) & 355M & 512 &  n/d & Español \\
Longformer-BSC & $\sim$130M & 4.096 & 0{,}416$^{b}$ & Clínico ES \\
IIC/RigoBERTa-2.0 & $\sim$600M & 512 & 0{,}74 & Español \\
IIC/RigoBERTa-Clinical & $\sim$600M & 512 & \textbf{0{,}799} & Clínico ES \\
mmBERT (JHU) & 307M & 8.192 & 0{,}12$^{b}$ & Multilingüe \\
ModernBERT & 149M & 8.192 & 0{,}601 & General \\
BETO & 110M & 512 &  n/d & Español \\
MiniLM multilingüe & 33M & 512 &  n/d & Embeddings \\
DistilBERT multilingüe & 66M & 512 &  n/d & Multilingüe \\
\hline
\end{tabular}

\smallskip\par\footnotesize Salvo indicación contraria, el F1 corresponde a la clasificación de capítulos (21 clases). $^{b}$~medido sobre los 809 códigos de bloque, tarea distinta y no comparable con el resto de la columna.
\end{table}

\section{Criterio de selección}

\texttt{IIC/RigoBERTa-Clinical} fue seleccionado para el entrenamiento en producción por tres razones: (1)~especialización en dominio clínico español, que garantiza representaciones adecuadas del vocabulario médico del corpus; (2)~mejor F1-micro experimental (0{,}799 en clasificación de capítulos) entre todos los modelos evaluados; y (3)~el mejor equilibrio entre F1-micro y F1-macro, lo que indica capacidad para clasificar tanto los códigos frecuentes como los minoritarios. \texttt{answerdotai/ModernBERT-base} queda identificado como candidato para iteraciones futuras gracias a su ventana de 8.192 tokens, aunque su F1-micro de 0{,}601 en nivel~1 refleja la penalización de operar sobre corpus en español con un modelo entrenado en inglés.
